\documentclass[11pt]{article}
\usepackage[margin=0.8in]{geometry}
\usepackage[T1]{fontenc}
\usepackage{hyperref}
\usepackage{longtable}
\usepackage{enumitem}

\title{Android F2FS FIEMAP GC Lab\\APK, ODEX, and VDEX Relocation Testing}
\author{}
\date{}

\begin{document}
\maketitle

\section{Goal}

The goal is to prove that Android package artifacts can keep the same content
and inode while their FIEMAP physical extents change due to F2FS garbage
collection.  The files of interest are normally:

\begin{itemize}[nosep]
\item \texttt{base.apk}
\item split \texttt{*.apk} files
\item \texttt{base.odex}
\item \texttt{base.vdex}
\end{itemize}

The desired outcome is not an application-level rewrite.  A successful GC
relocation result has these properties:

\begin{itemize}[nosep]
\item the SHA-256 hash before and after is identical;
\item the device and inode before and after are identical;
\item the FIEMAP physical extents differ;
\item F2FS GC counters, especially moved-block counters, increase.
\end{itemize}

This distinction matters.  Game play or settings changes usually mutate app
data, caches, shader caches, or save files.  They usually do not rewrite
\texttt{base.apk}, \texttt{base.odex}, or \texttt{base.vdex}.  For these files,
the interesting mechanism is F2FS moving still-valid blocks while reclaiming
partially invalid sections.

\section{Included Files}

\begin{longtable}{p{0.32\linewidth}p{0.60\linewidth}}
\textbf{File} & \textbf{Purpose} \\
\hline
\texttt{f2fs\_fiemap\_gc\_lab.py} &
Host-side Python harness.  It drives \texttt{adb}, discovers target files,
takes before/after snapshots, stages donor files, triggers GC, and writes a
summary. \\
\texttt{f2fs\_gc\_probe.c} &
Small native helper.  It issues \texttt{FS\_IOC\_FIEMAP},
\texttt{F2FS\_IOC\_GARBAGE\_COLLECT}, and
\texttt{F2FS\_IOC\_GARBAGE\_COLLECT\_RANGE}. \\
\texttt{Makefile} &
Build targets for a host compile check, Android helper build, and this PDF. \\
\texttt{README.md} &
Short command reference. \\
\end{longtable}

\section{Prerequisites}

\begin{itemize}
\item A rooted Android device.
\item Working \texttt{adb}.
\item \texttt{/data} mounted as F2FS.
\item Enough free space on \texttt{/data} for donor files.
\item Android NDK if you want the recommended native helper.
\end{itemize}

The helper is strongly recommended.  Many Android builds do not ship
\texttt{filefrag}, and there is no standard shell command for
\texttt{F2FS\_IOC\_GARBAGE\_COLLECT\_RANGE}.  Without the helper, the harness can
still collect package paths, hashes, stats, and F2FS sysfs counters, but FIEMAP
comparison may not be possible.

\section{Build the Helper}

From the kernel tree:

\begin{verbatim}
cd tools/f2fs_fiemap_gc_lab
export ANDROID_NDK_HOME=/path/to/android-ndk
make android ANDROID_ARCH=arm64 ANDROID_API=30
\end{verbatim}

For a quick local compile check only:

\begin{verbatim}
make host
\end{verbatim}

The Android build produces:

\begin{verbatim}
tools/f2fs_fiemap_gc_lab/f2fs_gc_probe.android
\end{verbatim}

The harness pushes this binary to \texttt{/data/local/tmp/f2fs\_gc\_probe} when
\texttt{--helper} is supplied.

\section{Recommended Workflow: Controlled APK Install}

This is the workflow to use when repeatability matters.

\subsection{Why this path is better}

F2FS segregates data by temperature and writes out of place.  Existing package
artifacts may be sitting in full, cold, valid sections.  GC has little reason to
move a fully valid section because it cannot reclaim useful space from it.

The controlled install workflow stages donor files with APK-like extensions,
keeps them alive during install, installs the target APK set, then removes the
donors before GC.  The intended effect is:

\begin{enumerate}[nosep]
\item donor files enter the same cold-data allocation class as APK-like files;
\item package artifacts are written while those donor allocations are present;
\item donor deletion creates invalid blocks near package artifacts;
\item GC now has a partially invalid section that may contain valid package
      blocks;
\item GC can move those valid package blocks and change FIEMAP.
\end{enumerate}

\subsection{Command}

For split APKs, repeat \texttt{--install-apk}.  For one APK, provide one
\texttt{--install-apk}.

\begin{verbatim}
python3 f2fs_fiemap_gc_lab.py \
  --package your.package.name \
  --helper ./f2fs_gc_probe.android \
  --install-apk /path/to/base.apk \
  --install-apk /path/to/split_config.apk \
  --range-gc \
  --gc-one-trials 128 \
  --gc-seconds 60 \
  --donor-count 64 \
  --donor-mib 4 \
  --collect-segment-info
\end{verbatim}

If you want to remove the existing package first:

\begin{verbatim}
python3 f2fs_fiemap_gc_lab.py \
  --package your.package.name \
  --helper ./f2fs_gc_probe.android \
  --install-apk /path/to/base.apk \
  --uninstall-first \
  --range-gc
\end{verbatim}

\section{Existing Play Store Install Workflow}

This is lower probability.  The harness cannot change how the package artifacts
were originally allocated.  Donor writes after installation may create general
GC pressure, but they may not dirty the sections containing the APK, ODEX, or
VDEX files.

\begin{verbatim}
python3 f2fs_fiemap_gc_lab.py \
  --package your.package.name \
  --helper ./f2fs_gc_probe.android \
  --range-gc \
  --gc-one-trials 128 \
  --gc-seconds 60 \
  --donor-count 64 \
  --donor-mib 4
\end{verbatim}

If package discovery misses an artifact, add it explicitly:

\begin{verbatim}
python3 f2fs_fiemap_gc_lab.py \
  --package your.package.name \
  --helper ./f2fs_gc_probe.android \
  --extra-path /data/app/.../oat/arm64/base.odex \
  --extra-path /data/app/.../oat/arm64/base.vdex \
  --range-gc
\end{verbatim}

\section{What the Harness Does}

\subsection{Step 1: Prove eligibility}

For every target file, the harness captures:

\begin{itemize}[nosep]
\item \texttt{stat} output;
\item SHA-256;
\item FIEMAP extents through the helper, or \texttt{filefrag} if available;
\item F2FS sysfs counters.
\end{itemize}

The result is stored under:

\begin{verbatim}
snapshots/before.json
snapshots/after.json
summary.txt
summary.json
\end{verbatim}

A file is considered relocated only when the content hash and inode remain
stable while the FIEMAP extent signature changes.

\subsection{Step 2: Manufacture partially invalid sections}

The donor-file strategy is controlled by:

\begin{verbatim}
--donor-dir /data/local/tmp/f2fs_gc_lab_donors
--donor-count 64
--donor-mib 4
--donor-exts apk,vdex,odex
\end{verbatim}

Each donor round creates one file per extension.  With the values above, the
temporary donor footprint is approximately:

\begin{verbatim}
64 rounds * 4 MiB * 3 extensions = 768 MiB
\end{verbatim}

Tune this for the free space on the device.  The default is smaller:

\begin{verbatim}
12 rounds * 8 MiB * 3 extensions = 288 MiB
\end{verbatim}

\subsection{Step 3: Force ART state before baseline}

Unless \texttt{--skip-dexopt} is supplied, the harness runs:

\begin{verbatim}
cmd package compile -m speed -f your.package.name
\end{verbatim}

This is done before the baseline snapshot so the test focuses on GC relocation,
not normal ART regeneration of ODEX or VDEX files.

\subsection{Step 4: Trigger GC}

The harness can use three GC paths.

\begin{itemize}
\item \texttt{--range-gc}: calls the helper's
      \texttt{F2FS\_IOC\_GARBAGE\_COLLECT\_RANGE} path for each target extent.
\item \texttt{--gc-one-trials N}: calls ordinary
      \texttt{F2FS\_IOC\_GARBAGE\_COLLECT} repeatedly through the helper.
\item \texttt{gc\_urgent}: writes F2FS sysfs knobs such as
      \texttt{gc\_urgent}, \texttt{gc\_urgent\_sleep\_time}, and
      \texttt{gc\_remaining\_trials}.
\end{itemize}

The relevant options are:

\begin{verbatim}
--range-gc
--gc-one-trials 128
--gc-trials 1000
--gc-seconds 60
--urgent-sleep-ms 20
\end{verbatim}

\subsection{Step 5: Compare and report}

The harness writes:

\begin{verbatim}
summary.txt
summary.json
logs/f2fs_attrs_before.txt
logs/f2fs_attrs_after.txt
logs/helper_gc_range.txt
logs/gc_urgent.txt
\end{verbatim}

A successful line in \texttt{summary.txt} looks like:

\begin{verbatim}
MOVED    sha=true  inode=true  extents=3->3 /data/app/.../base.apk
\end{verbatim}

That means content and inode stayed stable, but physical extents changed.

\section{Interpreting Common Outcomes}

\subsection{MOVED, hash true, inode true}

This is the desired result.  The file itself did not change, but FIEMAP changed.
This is consistent with F2FS relocating valid blocks during GC.

\subsection{MOVED, hash false}

This is not a clean GC relocation proof.  The file content changed.  For ODEX
and VDEX this may indicate ART regenerated the artifact.  Run dexopt before the
baseline snapshot, or use \texttt{--skip-dexopt} only when you intentionally
want to observe ART changes.

\subsection{same-map with GC counters increasing}

GC happened, but the target files were not moved.  Common reasons:

\begin{itemize}[nosep]
\item their sections were fully valid;
\item the sections were not selected as victims;
\item donor files landed in hot/warm data while the package artifacts were cold;
\item the files or sections were pinned;
\item there was not enough GC pressure;
\item \texttt{--range-gc} was not used.
\end{itemize}

\subsection{same-map with no moved-block counter delta}

The GC path probably did not move data blocks.  Check:

\begin{itemize}[nosep]
\item \texttt{/data} is really F2FS;
\item root shell has uid 0;
\item helper binary runs on the device architecture;
\item \texttt{gc\_urgent} and moved-block sysfs files exist;
\item the filesystem has enough dirty or prefree segments to make GC useful.
\end{itemize}

\section{Segment-Level Debug}

Use:

\begin{verbatim}
--collect-segment-info
\end{verbatim}

This collects:

\begin{verbatim}
logs/segment_info_before.txt
logs/segment_info_after.txt
logs/f2fs_status_before.txt
logs/f2fs_status_after.txt
\end{verbatim}

On kernels with the F2FS proc debug entries enabled, \texttt{segment\_info}
uses this format:

\begin{verbatim}
segment_type|valid_blocks
segment_type(0:HD, 1:WD, 2:CD, 3:HN, 4:WN, 5:CN)
\end{verbatim}

For a normal 4 KiB block, 2 MiB segment F2FS layout, one segment has 512
blocks.  If a target's section is full of valid blocks, GC has little reason to
move it unless a force-migration path is used.  The controlled install workflow
tries to make the target section partially invalid by deleting co-located donor
files before GC.

\section{F2FS Extension Classification}

F2FS can classify extensions as cold or hot through:

\begin{verbatim}
/sys/fs/f2fs/<device>/extension_list
\end{verbatim}

The harness records this in:

\begin{verbatim}
logs/f2fs_attrs_before.txt
logs/f2fs_attrs_after.txt
\end{verbatim}

If \texttt{apk}, \texttt{odex}, and \texttt{vdex} are cold extensions, donor
files with those suffixes are more likely to share the package artifact
allocation class.  If they are not cold extensions, the donor files may not help
as much.  Changing the extension list is a lab-only intervention because it
changes the filesystem policy being tested.

\section{Pinning Check}

F2FS supports pinned files.  Pinned files are intended not to be moved by GC
until enough GC failures cause F2FS to drop the guarantee.  The harness records:

\begin{verbatim}
gc_pin_file_thresh
\end{verbatim}

If target package artifacts are pinned, normal GC relocation may be very rare.
That is a filesystem policy result, not a harness failure.

\section{Safety Notes}

\begin{itemize}
\item Run this on a lab device or a disposable test image.
\item Keep enough free space on \texttt{/data}; donor files can be large.
\item The harness does not intentionally write target package files.
\item The harness does intentionally create and delete files under
      \texttt{/data/local/tmp/f2fs\_gc\_lab\_donors}.
\item Aggressive GC can affect device performance during the run.
\item If you use \texttt{--uninstall-first}, the package is removed before
      reinstall.
\end{itemize}

\section{Minimal Command Set}

Build helper:

\begin{verbatim}
cd tools/f2fs_fiemap_gc_lab
export ANDROID_NDK_HOME=/path/to/android-ndk
make android ANDROID_ARCH=arm64 ANDROID_API=30
\end{verbatim}

Run controlled install:

\begin{verbatim}
python3 f2fs_fiemap_gc_lab.py \
  --package your.package.name \
  --helper ./f2fs_gc_probe.android \
  --install-apk /path/to/base.apk \
  --range-gc \
  --gc-one-trials 128 \
  --gc-seconds 60 \
  --donor-count 64 \
  --donor-mib 4 \
  --collect-segment-info
\end{verbatim}

Inspect:

\begin{verbatim}
cat f2fs-gc-lab-*/summary.txt
\end{verbatim}

\section{Expected Reality}

A one-in-fifty result is plausible if the app was already installed and only
post-install app activity or unrelated random IO was used.  The package files
must be in victim sections that contain invalid blocks.  The controlled install
workflow is designed to create that condition.  If even the controlled flow does
not move \texttt{base.apk}, focus next on segment state, extension
classification, pinning, and whether the GC range ioctl is actually accepted by
the device kernel.

\end{document}
